% from ega4-2-fr.tex
\oldpage[IV-2]{177}
Soient $K$ un corps algébriquement clos de caractéristique 0, $B$ l'anneau
intègre
\[
K[U,V,W]/(U^2(U-W)-V^2(U+W))
\]
($U$, $V$, $W$ indéterminées) de sorte que $Y=\operatorname{Spec}(B)$ est
un « cône de sommet l'origine, ayant une génératrice double ». Nous
désignerons par $u$, $v$, $w$ les images de $U$, $V$, $W$ dans $B$. Soit
$R$ le corps des fractions de $B$, et considérons dans $R$ l'élément
$t=v(u+w)/u$, qui n'appartient pas à $B$; montrons que $C=B[t]$ est la
clôture intégrale de $B$. On a en effet $t^2=u^2-w^2$, donc $t$ est entier
sur $B$, et $v=tu/(u+w)$; l'anneau $C_1=K[t,u,w]$ est intégralement clos,
car il est isomorphe à
$K[T,U,W]/(T^2-U^2+W^2)$ et est donc la fermeture intégrale de l'anneau
intégralement clos $K[U,W]$ dans l'extension quadratique
$K(U,W)(\sqrt{U^2+W^2})$ de son corps des fractions (Bourbaki, Alg. comm.,
chap. V, § 1, n$^o$ 6, prop. 18). L'anneau de fractions
$K[t,u,w,1/(u+w)]$ de $C_1$ est donc aussi intégralement clos. De la même
manière, on voit que l'anneau $C_2=K[t,v,w]$ est intégralement clos, car
$t$ est racine de l'équation
$t(t-v)-w^2(2v-t)=0$, et par suite
$K[t,v,w,1/(t-v)]=K[t,tu/(u+w),w,(u+w)/(tw)]$ est intégralement clos. Finalement, compte tenu de ce
que $u$ et $w$ sont algébriquement indépendants sur $K$, on prouve aisément
que
\[
C=K[t,u,v,w]=K[t,u,w,1/(u+w)]\cap K[t,v,w,1/(t-v)],
\]
ce qui achève de montrer que $C$ est la clôture intégrale de $B$. Il est
immédiat que si $\mathfrak m_0$ est l'idéal maximal de $B$ engendré par
$u$, $v$, $w$ (« sommet du cône »), il existe un seul idéal maximal
$\mathfrak n_0$ de $C$ au-dessus de $\mathfrak m_0$, savoir l'idéal
engendré par $t$, $u$, $v$, $w$. Si l'on pose $A=B_{\mathfrak m_0}$, on en
déduit aisément que $A'=C_{\mathfrak n_0}$ est la clôture intégrale de
$A$, qui est donc unibranche, et par suite aussi géométriquement unibranche
puisque son corps résiduel $K$ est algébriquement clos. Mais dans $A$
l'idéal premier $\mathfrak p$ engendré par $u$ et $v$ est tel que la
clôture intégrale $A_\mathfrak p[t]$ de $A_\mathfrak p$ ne soit pas un
anneau local.

Nous verrons toutefois plus loin (9.7.10) que lorsque $X$ est un préschéma
localement noethérien tel que, si $X'$ est le normalisé de
$X_{\mathrm{red}}$, le morphisme canonique $X'\to X$ soit fini (ce qui sera
le cas si $X$ est tel que les anneaux de ses ouverts affines soient
universellement japonais (0, 23.1.1)), alors l'ensemble des points $x\in X$
où $X$ est géométriquement unibranche est localement constructible.

\begin{env}[6.15.3]
\label{IV.2.s6.15.3-fr}
Nous dirons pour abréger qu'un morphisme $f:X\to Y$ est \emph{radiciel en
un point} $y\in Y$, si $f^{-1}(y)$ est vide ou réduit à un seul point $x$
et si $k(x)$ est une extension radicielle de $k(y)$. Il revient donc au
même de dire que $f$ est radiciel en tous les points de $Y$ ou que $f$ est
radiciel (I, 3.5.8). Si
$g:f^{-1}(y)\to\operatorname{Spec}(k(y))$ est le morphisme image réciproque
de $f$ par le changement de base
$\operatorname{Spec}(k(y))\to Y$, il revient au même de dire que $f$ est
radiciel au point $y\in Y$, ou que $g$ est radiciel.
\end{env}

\begin{lemma}[6.15.3.1]
\label{IV.2.s6.15.3.1-fr}
\begin{enumerate}
\item[(i)] Soient $f:X\to Y$, $g:Y\to Z$ deux morphismes. Si $f$ est
radiciel en un point $y$ et $g$ radiciel au point $z=g(y)$, alors $g\circ f$
est radiciel au point $g(y)$. La réciproque est vraie si $f$ est surjectif.
\item[(ii)] Soient $f:X\to Y$, $h:Y'\to Y$ deux morphismes, et soit
$f'=f_{(Y')}:X\times_Y Y'\to Y'$. Soient $y'\in Y'$ et $y=h(y')$; alors
pour que $f$ soit radiciel au point $y$, il faut et il suffit que $f'$ soit
radiciel au point $y'$.
\end{enumerate}
\end{lemma}

\begin{proof}
